\ \\
Рассмотрим $\Z[i] = \{a + ib : a, b \in \mathbb{Z}\}$ -- \textit{гауссовы числа}.
\bigbreak
\noindent Покажем, что $\Z[i]^*$ = $\{\pm 1, \pm i\}$. Для этого определим на $\Z[i]$ норму $N : Z[i] \rightarrow \mathbb{N} \cup \{0\}$ следующим образом:
$ \forall a + ib \in \Z[i]$ положим $N(a + ib) := a^2 + b^2$ \\
Заметим, что тогда для всех $z, w \in \Z[i]$
выполнено $N(zw) = N(z)N(w)$ \\
\Proof $zw = z_1w_1 - w_2z_2 + i(w_1z_2 + w_2z_1) \;\; \Rightarrow \newline \Rightarrow\;\; N(zw) = (z_1w_1 - w_2z_2)^2 + (w_1z_2 + w_2z_1)^2 = (z_1^2 + z_2^2)(w_1^2 + w_2^2) = N(z)N(w)$ \;\; \EndProof
\\
Значит, все обратимые элементы в $\Z[i]$ имеют норму, равную единице, тогда непосредственной проверкой получим, что $\Z[i]^* = \{\pm 1, \pm i\}$.
\bigbreak
\noindent
На комплексной плоскости гауссовы числа образуют решетку следующего вида \newline (синим
отмечено $\Z[i]^*$, красным — пример четверки ассоциированных элементов):
\begin{center}
    \includegraphics[]{images/polina_1.PNG}
\end{center}